\section{Metodología}\label{sec:metodologia}

\textbf{Limpieza de las series.} Consideramos limpia una observación satelital con nubosidad de 20\,\% o menos y con
valor del índice. Elegimos ese umbral porque la nubosidad casi siempre vale 0 o 100\,\% (86\,\% de los registros) y
cuando vale 100 los índices vienen vacíos. Así cada parcela conserva de 19 a 34 observaciones de Sentinel-2 entre mayo
y octubre de 2025, con mediana de 25. Las observaciones descartadas no se imputan. Excluimos el índice de área foliar
(LAI) de Sentinel-2 porque su fórmula da valores negativos sobre suelo desnudo, sin significado físico, en 42\,\% de los
registros~\cite{ecuaciones}. Con las series limpias construimos 91 variables por parcela, como el máximo del NDVI de
2025, la fenología~\cite{zadoks1974}, la historia de 2022 a 2024, el clima y la topografía. Ninguna usa el rendimiento.

\textbf{Régimen.} Asignamos el régimen con una regla fija por nombre de municipio. No usamos el formato de cada
etiqueta, porque 7 de las 87 parcelas del régimen A tienen un solo decimal y esa regla las habría mandado al B. La regla
se conoce para las 59 parcelas, porque todas pertenecen a municipios con parcelas de entrenamiento. En seis parcelas de municipios con una a tres parcelas de entrenamiento el régimen se
dedujo de su propia etiqueta y de la cercanía a municipios del mismo formato. Sin ellas el RMSE pasa de 0.458 a
0.464\,t/ha, así que el efecto es pequeño.

\textbf{Esquema de evaluación.} Todo error se mide en parcelas que el modelo no vio. Usamos validación cruzada de 5
grupos, estratificada por municipio y repetida 10 veces, porque imita la partición oficial. Es el esquema A. Como la
validación al azar puede ser optimista cuando parcelas cercanas se parecen~\cite{roberts2017,ploton2020}, añadimos un
esquema B que deja fuera zonas completas, formadas al unir parcelas a menos de 2\,km. Todo lo que usa el rendimiento se
recalcula dentro de cada grupo con solo las parcelas de entrenamiento.

\textbf{Regla de decisión.} Con 138 parcelas una mejora puede aparecer por azar. Por eso una idea se acepta solo si baja
el RMSE en al menos 7 de las 10 repeticiones y si pasa una prueba de permutación con $p\le0.05$~\cite{ojala2010}. En la
prueba revolvemos las variables nuevas dentro de cada régimen y repetimos la validación 50 veces. Si el modelo predice
igual con ellas revueltas, la mejora no viene de esas variables.

\textbf{Rezago espacial.} Parcelas cercanas tienden a rendir parecido porque comparten suelo, clima y
manejo~\cite{tobler1970}. El rezago de una parcela es el promedio del rendimiento de hasta 10 parcelas de entrenamiento
del mismo régimen a 400\,m o menos, ponderado por el inverso de la distancia. Sin vecinos, recibe la media de su
régimen en el entrenamiento. Elegimos 400\,m porque a 250\,m solo el 25\,\% de las parcelas tiene vecino, mientras que
con radios mayores el promedio mezcla parcelas que ya no se parecen y el error sube.

\textbf{Algoritmo.} Comparamos modelos lineales, de árboles y un proceso gaussiano~\cite{rasmussen2006} con núcleo de
base radial y término de ruido. Sus tres parámetros se ajustan por máxima verosimilitud dentro de cada grupo.

\textbf{Intervalos.} Usamos predicción conformal de Mondrian~\cite{vovk2005,lei2018}, que calibra cada régimen por
separado porque sus errores se distribuyen distinto. A cada predicción le sumamos los cuantiles 0.10 y 0.90 de los
errores fuera de grupo de su régimen.

\textbf{Controles.} Probamos 73 variantes sin corregir por comparaciones múltiples, y elegir la mejor entre muchas
infla su resultado~\cite{cawley2010}. Lo corregimos con validación anidada, que repite la elección entre 12 variantes
dentro de cada grupo externo~\cite{varma2006}. Comparamos modelos con la prueba $t$ corregida de Nadeau y
Bengio~\cite{nadeau2003}, porque las repeticiones comparten parcelas. Una prueba automatizada verifica que alterar las
etiquetas de prueba no cambia las predicciones~\cite{kaufman2012}.

\textbf{Capa económica.} Valuamos la predicción con el precio medio rural del SIAP de 2023 a 2025~\cite{siap2026} y con
los costos del paquete de FIRA de primavera-verano 2022, el más reciente público, actualizados a 2025 con el índice
nacional de precios al consumidor~\cite{agrocostos}.
